\subsection{The complete kernel of the twelve-block Hurwitz action}

\begin{statement}{D23.F30}{Proved}{Twelve orthogonal simple factors and the exact binary quotient}
\end{statement}
Let $\Gamma=\langle A,B,C\rangle\leq\mathfrak S_{1344}$ be
the literal Hurwitz image of D23.P17. Let $Q_{12}$ be its action
on the twelve blocks recorded in
\path{agents/pentagon/hurwitz_cubic216_four_13_block0.npz},
and let $K$ be the kernel of that action. Use group-transversal
identifications of all blocks with block zero, as specified below.
The block-stabilizer image on that block is
\[
 L=\operatorname{PSO}_6^-(3),\qquad |L|=6531840,
 \qquad S=L'=\operatorname{P\Omega}_6^-(3)
              \simeq\operatorname{PSU}_4(3),
 \qquad |S|=3265920.
\]
Here the exact orthogonal identification and standard simplicity
input are D23.C23. Then
\[
 K'=S^{12},\qquad
 K/S^{12}=\mathcal C\subset\mathbb F_2^{12},\qquad
 \dim\mathcal C=9,
\]
where the twelve factors of $S^{12}$ act independently on the
twelve blocks. In the literal block labeling the four triples are
\[
 \mathcal T=\bigl\{
 \{0,9,11\},\{1,6,7\},\{2,4,10\},\{3,5,8\}\bigr\},
\]
and the code is exactly
\[
 \mathcal C=\left\{c\in\mathbb F_2^{12}:
           \sum_{i\in T}c_i\text{ is independent of }T\in\mathcal T
                              \right\}.
\]
More precisely, $K$ is the full preimage of this code in $L^{12}$.
The quotient $\Gamma/S^{12}$ has a faithful permutation
realization on $24$ points and order
\[
                         2^9\cdot216=110592.
\]
The socle and the eventual derived subgroup of $\Gamma$ are
both $S^{12}$. Thus the displayed order-$110592$ group is its
largest solvable quotient.

\begin{proof}
Write $q_A,q_B,q_C\in\mathfrak S_{12}$ for the three block
permutations and put $I_0=\{v:\operatorname{block}(v)=0\}$.
A breadth-first traversal on the twelve blocks chooses
$T_j\in\Gamma$ with $T_j(I_0)=I_j$ and $T_0=1$. If a
generator $g$ carries $i$ to $j$, the exact transition
\[
                    T_j^{-1}gT_i
\]
stabilizes $I_0$. The Schreier lemma therefore says that these
transitions generate the entire block stabilizer, before and
after restriction to $I_0$. The certificate retains the literal
braid words and all the $1344$-point permutations. Their
restrictions give thirteen distinct nonidentity generators,
which reduce to four generators without changing the group.
Exact Schreier--Sims gives order $6531840$, derived order
$3265920$, and a perfect derived group. The stabilizer of one
of the $112$ points has orbit sizes $1,30,81$. D23.C23
identifies these same generators as the full group $L$ above;
that identification is not inferred merely from their order.

Similarly enumerate all $216$ elements of $Q_{12}$ and choose
word lifts $T_q\in\Gamma$ of them. For each generator $g$ and
each $q\in Q_{12}$, the word
\[
                       T_{q_gq}^{-1}gT_q
\]
lies in $K$. These words generate $K$ by the Schreier lemma.
There are $153$ distinct nonidentity full permutations in this
Schreier list. Distinct full permutations are retained even when
their block-zero restrictions agree or are the identity.
Their block-zero image has order $6531840$, the same as $L$, so the
projection of $K$ onto that block is all of $L$. Conjugation
by the $T_i$ proves this for every block. Since $K$ fixes all
blocks, its action identifies it with a subgroup of $L^{12}$.
The action is faithful because the blocks cover all $1344$
points.

We next check the correlations, not just the separate projections.
The stabilizer of zero in $Q_{12}$ has precisely the orbits
\[
 \{0\},\quad\{9,11\},\quad\{1,2,3,4,5,6,7,8,10\}.
\]
There are therefore only two ordered-pair types for distinct
blocks. Normalize the restriction of each of the above
Schreier generators to block $i$ by conjugating with $T_i$.
The joint restrictions to $(0,9)$ and to $(0,1)$ both generate
groups of order
\[
          42664933785600=6531840^2.
\]
Each lies in $L\times L$ and projects onto both factors;
equality of orders proves that each is the full product.
By the two orbit types, the projection of $K$ onto any two
distinct blocks is therefore $L\times L$. Its derived
projection is $S\times S$, of exact order
$10666233446400=3265920^2$.

For clarity, the simple-product inference used next needs
nonabelian simplicity. If a subgroup $J\leq S^m$ projects
onto every factor and every pair, where $S$ is nonabelian
simple, then $J=S^m$. Prove this by induction. Having obtained
the first $m-1$ factors, Goursat's lemma between $S^{m-1}$
and the last factor either gives the full product or identifies
a simple quotient of $S^{m-1}$ with the last $S$. Every normal
subgroup of $S^{m-1}$ is a product of factors: commuting an
element of such a normal subgroup with one factor isolates
that coordinate, and simplicity and trivial center force
each nontrivial projected factor to be contained in the
normal subgroup. Thus a simple quotient selects one factor.
The resulting diagonal relation contradicts the assumed full
projection onto that pair. This proves the induction.
Apply the lemma to $K'\leq S^{12}$. Its single and pair
projections are $S$ and $S\times S$, respectively, so
$K'=S^{12}$.

The coordinate characters $L\to L/S=C_2$ consequently
identify $K/S^{12}$ with a binary linear code. On the $153$
Schreier generators the exact computation tests membership
in the derived subgroup of $L$ in each coordinate. Gaussian
elimination on the resulting $153$ by $12$ matrix gives rank
nine. Every row has the same parity on the four triples
displayed in the statement. Those triples are exactly the
images of $\{0,9,11\}$ under $Q_{12}$. The equality of their
four parities imposes three independent equations, hence
defines a nine-dimensional subspace. The inclusion and the
rank computation prove the displayed equality for $\mathcal C$.
As $K$ already contains $S^{12}$, it is the full preimage of
that code in $L^{12}$.

Here is the faithful quotient explicitly. For $g\in\Gamma$
and a block $i$, put $j=q_g(i)$ and let
\[
  e_g(i)\in\mathbb F_2
       \quad\text{be the image of }T_j^{-1}gT_i\text{ in }L/S.
\]
On the $24$ pairs $(i,e)\in\{0,\ldots,11\}\times\mathbb F_2$
define
\[
                    g(i,e)=(q_g(i),e+e_g(i)).
\]
The transition identity proves multiplicativity. Its kernel
consists precisely of those elements of $K$ whose twelve
coordinates belong to $S$, namely $S^{12}$. Its order is
$|\mathcal C|\,|Q_{12}|=2^9\cdot216$, also independently
checked directly from the three resulting permutations.
The block quotient has derived orders
$216,72,18,9,1$, so this quotient is solvable. Since $S^{12}$
is perfect, it belongs to every derived subgroup of $\Gamma$;
solvability of the quotient proves the eventual-derived claim.

Finally $S^{12}$ is a minimal normal subgroup of $\Gamma$:
normal subgroups of this direct product are products of its
simple factors, and $\Gamma$ acts transitively on them.
Its centralizer in the full permutation group is trivial.
Indeed a commuting permutation must preserve each factor's
support, and the centralizer of $S$ on one block is trivial:
a permutation commuting with a transitive group is determined
by the image of a point, which must be fixed by its point
stabilizer, whose orbit sizes $1,30,81$ leave just that point.
Any different minimal normal subgroup of $\Gamma$ would
centralize $S^{12}$, which is impossible. This proves the
socle assertion.
\end{proof}

\paragraph{Certificates and inputs.}
The complete Schreier construction and exact order checks are
\path{agents/five_level_six/hurwitz_block_stabilizer.py}
and its JSON/NPZ records. The joint restrictions, all local
transports, binary Gaussian elimination, signed quotient,
and the directly verified orbital graph parameters
$(112,30,2,10)$ are in
\path{agents/five_level_six/hurwitz_block_correlations.py}
and its JSON/NPZ records. The classical-group identification
and simplicity are explicit inputs from D23.C23. No abstract
group name or simple-product conclusion is obtained from
coinciding orders alone.

\subsection{The solvable block quotient is the earlier Burau quotient}

\begin{statement}{D23.F31}{Proved}{An exact marked Burau identification}
\end{statement}
The order-$110592$ signed permutation quotient in D23.F30 is
marked-isomorphic to the group generated by the literal
matrices of supplied FP.36.2 over
$R=(\mathbb Z/4)[\omega]/(\omega^2+\omega+1)$, modulo
$\langle\omega I\rangle$:
\[
 A_B=\begin{pmatrix}1-\omega&\omega&0\\1&0&0\\0&0&1\end{pmatrix},
 \quad B_B=\begin{pmatrix}1&0&0\\0&1-\omega&\omega\\0&1&0\end{pmatrix},
 \quad C_B=\begin{pmatrix}1&0&-1\\0&1&-1\\0&0&-\omega\end{pmatrix}.
\]
The isomorphism sends the positive Hurwitz generators
$A,B,C$ to exactly $A_B,B_B,C_B$. The simultaneous flip of
all twelve binary pairs maps to $-I$. The projective quotient
therefore has order $55296$ and is the already supplied
Burau quotient, with its original markings.

\begin{proof}
Encode $a+b\omega\in R$ by $a+4b$ with $0\leq a,b<4$.
The certificate uses the exact multiplication
\[
 (a+b\omega)(c+d\omega)
     =(ac-bd)+(ad+bc-bd)\omega
\]
modulo four. It constructs the three matrices displayed
above and identifies a matrix with its three scalar multiples
by $1,\omega,\omega^2$, choosing a canonical representative.

Starting with the identity, simultaneously traverse the signed
permutation group and these matrix classes using the three
marked positive generators. Whenever a permutation already
occurred, check that its newly computed matrix class is
exactly the stored one. This check is made for every one of
the $3\cdot110592$ directed generator edges. All checks pass.
There are $110592$ permutations and also $110592$ distinct
matrix classes. Thus the marked homomorphism is well defined
and injective as well as surjective. In particular this is an
explicit marked isomorphism, not merely a match of orders or
derived series. The code also finds the global binary flip
and checks that its matrix class is precisely $-I$.

The determinants of all three matrix generators are $-\omega$.
Their determinant values therefore lie in $\mu_6\subset R^\times$,
and this determinant character is unchanged by multiplying
the matrix by $\omega I$. The character
\[
                 \epsilon(g)=\det(g)^3\in\{1,-1\}
\]
takes $-I$ to $-1$. Hence the central scalar subgroup
$\langle-I\rangle$ splits off, and projection to the
projective group restricts to an isomorphism on
$\ker\epsilon$. Its order is $110592/2=55296$.
This proves the projective identification with the exact
matrices and central relation of FP.36.
\end{proof}

\paragraph{Consequences and exact limitation.}
The signed quotient has first-pair image of order $192$.
Adjoining the marked cubic detector of D23.P10 raises the
joint image order to $3456$, so that detector is not a marked
quotient of this first-pair image. Its mixed image has order
$144$, and adjoining the marked $S_4$ generators
$(123),(1234)$ leaves the order $144$, proving retention of
that harmonic word test. These are checked in the same
correlation certificate.

For every commutator word $f$ evaluated at any two elements of
the signed quotient, $\epsilon(f)=1$. Thus its projective
value determines its entire signed value; the extra scalar
$-I$ supplies no new commutator-word constraint. In particular
the signed quotient alone adds no new identity comparing
the cubic frame with the harmonic or half-Kummer character.
This does not assert a general failure of such an identity on
actual gentle elements. The new possible information in the
$1344$-point construction lies in its twelve orthogonal
factors and their actual full-GT normalizer, not in this
already known solvable quotient.

The exhaustive marked comparison is
\path{agents/five_level_six/hurwitz_signed_burau.py}; its
JSON records all counts and the scalar-flip check. The proof
uses the matrices and markings supplied in FP.36, and gives
a new marked identification of the Hurwitz quotient with them.
